% This section answers RQ2:
% \textit{What compiler-level and architectural factors explain
% the underperformance of DSL-written kernels on convolution
% and other affected kernel categories?}

% \noindent\textit{Evidence basis:}
% Root causes RC0--RC4 are motivated by the latency measurements in \cref{sec:evaluation}
% and community reports~\cite{triton2022conv591}.
% Hardware performance counter confirmation via Nsight Compute (\cref{sec:meth:profiling})
% is pending; claims marked \textbf{[counter pending]} remain hypotheses awaiting profiling.

% \subsection{RC0: TileLang Compiler-Level Deficiencies}
% \label{sec:analysis:jit}

% \textbf{Finding: TileLang's compiled kernels exhibit two structural performance deficiencies:
% (i) use of \texttt{T.serial} manual reduction loops instead of the native parallel \texttt{T.reduce} primitive,
% executing the reduction sequentially over each partial value; and
% (ii) absent 128-bit vectorized memory loads for normalization kernels.
% These compound the per-kernel performance gaps measured in RQ1.
% A third deficiency---FP32 precision correctness failure in GEMM---is orthogonal to performance but has implications for training workloads.}

% Across all normalization and reduction benchmarks,
% TileLang reports anomalously high latencies regardless of kernel type:
% for large inputs, LayerNorm takes 273~ms (vs.\ Triton 0.92~ms and PyTorch 0.87~ms),
% and max reduction takes 28.1~ms (vs.\ PyTorch 1.62~ms).
% This deficit is qualitatively distinct from RC1--RC4 and compounds them.

% Two mechanisms contribute to RC0.
% First, the original TileLang normalization kernels implement the reduction over the normalized dimension using \texttt{T.serial} loops---TileLang's sequential for-loop construct---iterating over each partial sum one at a time.
% For LayerNorm, which accumulates 8 bfloat16 partial statistics per thread (sum and sum-of-squares for the mean and variance passes), this means the reduction is computed in 8 sequential steps with no inter-thread parallelism.
% The native alternative, \texttt{T.reduce}, lowers to a parallel butterfly reduction via \texttt{tl::AllReduce}.
% A \emph{butterfly reduction} is a logarithmic communication pattern in which threads exchange partial results in $\log_2 N$ rounds (where $N$ is the number of participating threads);
% over 256 threads grouped into warps of 32, each round requires one inter-warp synchronization barrier, for $\log_2(256/32) = 3$ total barriers---applied once for all 8 values simultaneously.
% Replacing the \texttt{T.serial} loop with a single \texttt{T.reduce} call reduces the dominant latency source: in our experiments, LayerNorm drops from 1090~ms to 5.24~ms ($208\times$ improvement, \cref{sec:mitig:norm}).

% Second, TileLang's compiled normalization kernels do not issue 128-bit vectorized loads
% (\texttt{LDG.128}),
% instead fetching 16-bit bfloat16 elements individually.
% For a tensor of shape $8192 \times 8192$,
% this produces an excessive number of discrete memory transactions
% that bottlenecks the memory controller
% and prevents the L2 cache from streaming coalesced data to the SMs.

% \paragraph{TileLang float32 correctness.}
% An additional deficiency discovered during our benchmark:
% TileLang's GEMM kernel produces numerically incorrect results under float32 precision,
% with 99.6\% of output elements mismatched against the PyTorch reference
% (greatest relative difference: $2067\times$ at tested problem shapes of $4096 \times 2048 \times 1024$).
% The float16 variant passes correctness checks in all cases.
% For this reason, all TileLang GEMM measurements in this paper use FP16;
% float32 TileLang GEMM results are excluded.
% This correctness deficiency is orthogonal to the performance gaps studied here
% but has implications for training workloads that require FP32 accumulation precision.

% Both deficiencies are correctable within the existing TileLang API:
% (i) replacing \texttt{T.serial} reduction loops with the native \texttt{T.reduce} primitive (already available and used correctly elsewhere in the benchmark suite), and
% (ii) enabling 128-bit vectorized global-memory loads for normalization kernels.

% \subsection{RC1: Absent Vectorization of Strided Memory Accesses}
% \label{sec:analysis:vec}

% \textbf{Finding: Triton-generated convolution code fails to issue vectorized (128-bit) load instructions,
% preventing asynchronous copy (\texttt{cp.async}) from being applied and collapsing the software pipeline.}

% NVIDIA Ampere's memory subsystem achieves peak bandwidth only when global memory loads
% are issued as 128-bit vector instructions (\texttt{LDG.128}).
% On Hopper, the Tensor Memory Accelerator (TMA) similarly requires
% aligned, contiguous access descriptors.
% For GEMM kernels, Triton's compiler successfully emits vectorized loads
% because the data layout is row-major and tile boundaries are aligned to 16-byte boundaries.

% Convolution complicates this in two ways.
% First, each output element requires gathering input values from a
% $(K_H \times K_W)$-window strided by \texttt{(stride\_h, stride\_w)} in the spatial dimensions.
% The resulting access pattern is non-contiguous across the innermost dimension
% when the input is stored in NHWC layout,
% breaking the alignment requirements for \texttt{LDG.128}.
% Second, for filter sizes $K_H \times K_W > 1$,
% the compiler must prove at compile time that successive iterations of the
% spatial reduction loop produce aligned addresses;
% this proof is not discharged by Triton's current alias analysis,
% so it conservatively emits scalar 32-bit loads.

% The consequence is a cascade noted in the community~\cite{triton2022conv591}:
% un-vectorized accesses prevent \texttt{cp.async} from firing,
% because CUDA's asynchronous copy instructions require the source address
% to produce a 128-bit aligned transfer.
% Without \texttt{cp.async}, the software pipeline cannot overlap data movement with computation,
% and the kernel stalls on global memory latency at every tile boundary.

% \paragraph{Vectorized loads on the RTX 4000 Ada.}
% The same deficiency manifests with added severity on Ada Lovelace hardware,
% where the 160-bit GDDR6 memory interface (360~GB/s peak bandwidth)
% requires 128-bit vector loads (\texttt{LDG.128}, fetching eight fp16 values simultaneously)
% to saturate the memory bus.
% For a convolution kernel that falls back to scalar 32-bit loads,
% the effective achievable bandwidth is one-quarter of peak.
% PyTorch's cuDNN implementation explicitly benchmarks and selects
% the fastest algorithm for each convolution configuration,
% with its NHWC implicit GEMM path issuing vectorized loads unconditionally;
% Triton and TileLang currently lack equivalent dispatching logic for convolution.

% % TODO: insert Nsight Compute counter table showing vectorized-load fraction,
% %       cp.async issue count, and warp stall cycles for GEMM vs. Conv2d

% \subsection{RC2: Auto-Tuning Search Space Mismatch}
% \label{sec:analysis:autotune}

% \textbf{Finding: The static tile-configuration search space in Triton's \texttt{@triton.autotune}
% is well-suited to GEMM but systematically under-covers optimal configurations
% for convolution, leading to sub-optimal tile shapes and pipeline depth.}

% Triton's auto-tuner sweeps a programmer-specified list of configurations
% $\{(\texttt{BLOCK\_M}, \texttt{BLOCK\_N}, \texttt{BLOCK\_K}, \texttt{num\_stages}, \texttt{num\_warps})\}$.
% For GEMM, this 5-dimensional space is well-understood:
% large square tiles ($128 \times 128$) with 4--8 pipeline stages
% are optimal across most A100 problem shapes~\cite{tillet2019triton}.
% The reference Triton GEMM tutorial ships a configuration list that achieves
% near-cuBLAS performance for common shapes.

% This is further evidenced by the large-matmul result in \cref{sec:eval:gemm}:
% Triton achieves only 31.7\% of cuBLAS throughput on a $16384^2$ FP16 matrix multiplication
% (361.9~ms vs.\ PyTorch's 114.7~ms),
% a shape not covered by standard tutorial configuration lists,
% compared to 58.2\% on a $128\times2048^2$ batched GEMM
% where the auto-tuner has well-populated configurations.
% The convolution gap (34.9\% for the $32\times256\times128\times128$, $3\times3$ configuration)
% follows the same pattern at an even larger scale.

% \paragraph{L2 cache residency and the GEMM divide.}
% For a $16384^2$ FP16 matrix multiplication,
% the combined memory footprint of matrices A, B, and C is approximately
% $3 \times 16384^2 \times 2 \approx 1.6$~GB,
% vastly exceeding the 48~MB L2 cache capacity of the RTX 4000 Ada.
% At this scale, performance is entirely determined by how effectively
% tiling strategies keep working data resident in L2 across tile iterations.
% cuBLAS employs hardware-specific heuristics (via cuBLASLt's plan selector)
% to slice the matrix sequence such that chunks remain L2-resident for as long as possible;
% Triton's generic \texttt{tl.dot} compilation lacks equivalent cache-blocking heuristics,
% resulting in repeated eviction and re-fetching of the same data
% and the 3.2$\times$ penalty observed in the benchmarks
% (114.7~ms for PyTorch vs.\ 361.9~ms for Triton on $16384^2$).
% TileLang (56.4\%) partially avoids this because its explicit schedule
% specifies pipeline stages and memory staging more directly.

% Convolution introduces two additional degrees of freedom:
% filter size $(K_H, K_W)$ and the reduction over the spatial window.
% Optimal tile shapes depend strongly on filter size ---
% $1 \times 1$ convolutions behave like GEMM (large tiles preferred),
% while $3 \times 3$ convolutions require smaller $K$-dimension tiles
% to keep shared memory within 48 KB per tile
% (the L1 capacity of a single SM on A100~\cite{cublas2024}).
% The default configuration lists shipped with TritonBench's convolution kernels
% do not include configurations specialized for $5 \times 5$ and $7 \times 7$ filters,
% leaving substantial performance on the table.

% Furthermore, TileLang's explicit pipeline declarations
% require the programmer to specify \texttt{num\_stages} ahead of time.
% For convolutions with large filters, the optimal number of pipeline stages
% depends on the register pressure from accumulating partial sums
% across the spatial loop --- a dependency that current scheduling heuristics do not model.

% % TODO: quantify via extended search-space sweep

% \subsection{RC3: Register Pressure and Occupancy Reduction}
% \label{sec:analysis:regpressure}

% \textbf{Finding: Convolution kernels with $5 \times 5$ and larger filters
% exceed the per-thread register budget, causing spill to local memory
% and reducing SM occupancy below the level required for latency hiding.}

% Each streaming multiprocessor on A100 has 65,536 32-bit registers per block.
% A block of 128 threads with 64 registers per thread occupies the full register file,
% preventing a second block from being resident and eliminating pipeline interleaving.
% For large-filter convolutions, maintaining partial sums for the
% $K_H \times K_W$ reduction --- plus input tile and output tile registers ---
% pushes per-thread register usage well above 64 registers.
% The Triton compiler cannot expose register count as a tunable parameter;
% \texttt{num\_warps} controls the thread count per block but does not
% directly constrain per-thread register usage,
% leaving the compiler to make potentially conservative register allocation decisions.

% % TODO: insert table of per-thread register count vs. filter size, measured via cuobjdump

% \subsection{RC4: Absence of Algorithm-Level Diversity}
% \label{sec:analysis:algo}

% \textbf{Finding: Both Triton and TileLang implement only the im2col-GEMM lowering for convolution,
% providing no access to Winograd or FFT algorithms that cuDNN uses for specific filter sizes.}

% cuDNN's runtime selector chooses Winograd for $3 \times 3$ filters
% when arithmetic reduction is the bottleneck.
% Winograd convolution reduces the multiply-accumulate count by a factor of
% approximately $2.25\times$ for $3 \times 3$ (Winograd $F(2,3)$) filters,
% at the cost of additional data transformation passes.
% Neither Triton nor TileLang expose the Winograd transformation
% as a schedulable primitive,
% so they cannot exploit this arithmetic reduction.
% This contributes to the gap specifically on $3 \times 3$ filters
% when the computation is arithmetic-bound.

% % TODO: estimate theoretical impact of Winograd for each tested configuration

% \subsection{Summary of Root Causes}
% \label{sec:analysis:summary}

% \Cref{tab:rootcauses} maps each root cause to the kernel category it primarily affects,
% its estimated contribution to the throughput gap, and the Nsight Compute counter
% most diagnostic for its detection.

% \begin{table*}[t]
%   \centering
%   \caption{Root-cause summary. Measured contribution is the latency speedup recovered by the corresponding mitigation (\cref{sec:mitigation}); ``---'' indicates no direct mitigation experiment conducted. RC0a and RC0b are TileLang-specific; RC1--RC4 affect both DSLs.}
%   \label{tab:rootcauses}
%   \begin{tabular}{p{5cm}p{3.5cm}p{2.2cm}p{4.5cm}}
%     \toprule
%     Root cause                     & Affected kernels              & Contrib. & Diagnostic counter              \\
%     \midrule
%     RC0a: \texttt{T.serial} instead of \texttt{T.reduce} & All TileLang reductions, norm & $208\times$ (LayerNorm) & Thread sync count, warp stall   \\
%     RC0b: Absent vectorized loads / dtype cast            & TileLang norm, conv           & $3.4\times$ (LayerNorm) & \texttt{l1tex\_\_t\_bytes}      \\
%     RC1: Strided access, scalar LD  & Conv2d ($K > 1$), Triton      & $2.57\times$ with RC2   & Vectorized-load fraction        \\
%     RC2a: Auto-tuning mismatch      & Matmul, Conv2d, Depthwise     & $1.66\times$ (Matmul)   & TFLOPS vs.\ swept configs       \\
%     RC2b: L2 cache residency        & Matmul $\geq 16384^2$         & ---                     & L2 hit rate, DRAM throughput    \\
%     RC3: Register pressure          & Conv2d ($K \geq 5$)           & ---                     & \texttt{sm\_\_register\_spill} \\
%     RC4: No Winograd                & Conv2d $3\times3$             & ---                     & SM utilization delta            \\
%     \bottomrule
%   \end{tabular}
% \end{table*}

\noindent\textit{The following root causes are motivated by latency measurements (\cref{sec:evaluation}) and community reports~\cite{triton2022conv591}.}

\phantomsection\label{sec:analysis:jit}
\paragraph{RC0: TileLang Compiler-Level Deficiencies.}
% REVISION-NOTE-START id=4134A-W1 severity=Major
\textcolor{red}{\textbf{[REVISION \#4134A-W1+GAP-4, Major]}
Concern (\textit{R1}): The submitted-PDF label conflates kernel-authoring (T.serial choice) with compiler-codegen (LDG.128 absent); the rebuttal splits them as RC0a (authoring, user-space-fixable) and RC0b (codegen, compiler-fixable).
Evidence: \texttt{../REVISION\_TODO.md}, \texttt{../experiments/results/NVIDIA\_RTX\_4000\_Ada\_Generation/NCU\_FINDINGS.md}.
Fix: Restructure RC0 as two subheads: ``RC0a: TileLang Reduction Authoring (T.serial $\to$ T.reduce)'' for the user-space-fixable mechanism and ``RC0b: TileLang Vectorized-Load Codegen'' for the compiler-fixable one, keeping the body content but splitting the framing.}
% REVISION-NOTE-END id=4134A-W1


Two mechanisms contribute to RC0.
First, the original TileLang normalization kernels implement the reduction over the normalized dimension using \texttt{T.serial} loops---TileLang's sequential for-loop construct---iterating over each partial sum one at a time.
For LayerNorm, which accumulates 8 bfloat16 partial statistics per thread (sum and sum-of-squares for the mean and variance passes), this means the reduction is computed in 8 sequential steps with no inter-thread parallelism.
The native alternative, \texttt{T.reduce}, lowers to a parallel butterfly reduction via \texttt{tl::AllReduce}.
A \emph{butterfly reduction} is a logarithmic communication pattern in which threads exchange partial results in $\log_2 N$ rounds ($\log_2(256/32) = 3$ barriers over 256 threads grouped into warps of 32).
% REVISION-NOTE-START id=4134A-W1 severity=Major
\textcolor{red}{\textbf{[REVISION \#4134A-W1+GAP-1, Major]}
Concern (\textit{R1}): This sentence presents barriers as the dominant cost, but NCU shows \texttt{warp\_stall\_barrier} $\approx$ 0; the dominant stall is \texttt{warp\_stall\_long\_scoreboard} $\approx$ 105 cycles (memory latency).
Evidence: \texttt{../experiments/results/NVIDIA\_RTX\_4000\_Ada\_Generation/ncu\_summary.csv}.
Fix: Keep butterfly-reduction description as algorithmic context, but add: ``The performance gain from \texttt{T.reduce} comes from better memory-latency hiding (\texttt{warp\_stall\_long\_scoreboard} drops from 104.91 to $<$ 5 cycles), not from removing barriers (barrier stalls were already $\approx$ 0 in both kernels).''}
% REVISION-NOTE-END id=4134A-W1
Replacing the \texttt{T.serial} loop with a single \texttt{T.reduce} call reduces the dominant latency source: in our experiments, LayerNorm drops from 1090~ms to 5.24~ms ($208\times$ improvement, \cref{sec:mitig:norm}).

Second, TileLang's compiled normalization kernels do not issue 128-bit vectorized loads
(\texttt{LDG.128}),
instead fetching 16-bit bfloat16 elements individually.
For a tensor of shape $8192 \times 8192$,
this produces an excessive number of discrete memory transactions
that bottlenecks the memory controller
and prevents the L2 cache from streaming coalesced data to the SMs.

% REVISION-NOTE-START id=4134A-W1 severity=Major
\textcolor{red}{\textbf{[REVISION \#4134A-W1+GAP-1+GAP-4, Major]}
Concern (\textit{R1}): This is the RC0b mechanism; label it explicitly.
Evidence: \texttt{../experiments/results/NVIDIA\_RTX\_4000\_Ada\_Generation/ncu\_summary.csv}.
Fix: Re-label this paragraph as ``RC0b: Absent Vectorized Loads (Compiler Codegen)'' and tighten the wording so it does not blur with RC0a.}
% REVISION-NOTE-END id=4134A-W1

\paragraph{TileLang float32 correctness.}
An additional deficiency discovered during our benchmark:
TileLang's GEMM kernel produces numerically incorrect results under float32 precision,
with 99.6\% of output elements mismatched against the PyTorch reference
(greatest relative difference: $2067\times$ at tested problem shapes of $4096 \times 2048 \times 1024$).
The float16 variant passes correctness checks in all cases.
For this reason, all TileLang GEMM measurements in this paper use FP16;
float32 TileLang GEMM results are excluded.
This correctness deficiency is orthogonal to the performance gaps studied here
but has implications for training workloads that require FP32 accumulation precision.

% REVISION-NOTE-START id=4134A-W1 severity=Major
\textcolor{red}{\textbf{[REVISION \#4134A-W1+4134A-W11, Major]}
Concern (\textit{R1}): Currently flagged as orthogonal; root cause is silent TF32 lowering in \texttt{T.gemm}, not a logic error.
Evidence: \texttt{../experiments/results/NVIDIA\_RTX\_4000\_Ada\_Generation/fp32\_gemm.csv}.
Fix: Reframe the 99.6\% mismatch as consistent with TF32 mantissa truncation (10-bit mantissa vs FP32's 23-bit), reproducible across four isolation arms in \texttt{exp\_fp32\_gemm.py}; cite the experiment and add a forward ref to a TF32-lowering discussion.}
% REVISION-NOTE-END id=4134A-W1


\textbf{Finding RC0: TileLang's compiled kernels exhibit two structural performance deficiencies:
(i) use of \texttt{T.serial} manual reduction loops instead of the native parallel \texttt{T.reduce} primitive,
executing the reduction sequentially over each partial value; and
(ii) absent 128-bit vectorized memory loads for normalization kernels.
These compound the per-kernel performance gaps measured in RQ1.
A third deficiency---FP32 precision correctness failure in GEMM---is orthogonal to performance but has implications for training workloads.}

% REVISION-NOTE-START id=4134A-W1 severity=Major
\textcolor{red}{\textbf{[REVISION \#4134A-W1+GAP-1+GAP-4, Major]}
Concern (\textit{R1}): The bold finding should reflect the RC0a/RC0b split and the memory-latency framing.
Evidence: \texttt{../experiments/results/NVIDIA\_RTX\_4000\_Ada\_Generation/NCU\_FINDINGS.md}.
Fix: Rewrite the finding: ``TileLang compiled kernels exhibit two structural deficiencies: (RC0a) \texttt{T.serial} reductions serialize memory-latency hiding rather than synchronization (\texttt{warp\_stall\_long\_scoreboard} dominates; \texttt{barrier} $\approx$ 0), correctable in user space via \texttt{T.reduce}; and (RC0b) absent 128-bit vectorized memory loads for normalization, which require compiler fixes. A third deficiency --- silent TF32 lowering of \texttt{T.gemm} under FP32 --- is correctness-affecting but orthogonal to performance.''}
% REVISION-NOTE-END id=4134A-W1

\phantomsection\label{sec:analysis:vec}
\paragraph{RC1: Absent Vectorization of Strided Memory Accesses.}
NVIDIA Ampere's memory subsystem achieves peak bandwidth only when global memory loads
are issued as 128-bit vector instructions (\texttt{LDG.128}).
On Hopper, the Tensor Memory Accelerator (TMA) similarly requires
aligned, contiguous access descriptors.
For GEMM kernels, Triton's compiler successfully emits vectorized loads
because the data layout is row-major and tile boundaries are aligned to 16-byte boundaries.

Convolution complicates this in two ways.
First, each output element requires gathering input values from a
$(K_H \times K_W)$-window strided by \texttt{(stride\_h, stride\_w)} in the spatial dimensions.
The resulting access pattern is non-contiguous across the innermost dimension
when the input is stored in NHWC layout,
breaking the alignment requirements for \texttt{LDG.128}.
Second, for filter sizes $K_H \times K_W > 1$,
the compiler must prove at compile time that successive iterations of the
spatial reduction loop produce aligned addresses;
this proof is not discharged by Triton's current alias analysis,
so it conservatively emits scalar 32-bit loads.

The consequence is a cascade noted in the community~\cite{triton2022conv591}:
un-vectorized accesses prevent \texttt{cp.async} from firing,
because CUDA's asynchronous copy instructions require the source address
to produce a 128-bit aligned transfer.
Without \texttt{cp.async}, the software pipeline cannot overlap data movement with computation,
and the kernel stalls on global memory latency at every tile boundary.
On the RTX 4000 Ada's 160-bit GDDR6 interface (360~GB/s peak), scalar 32-bit loads achieve at most one-quarter of peak bandwidth; \texttt{LDG.128} vectorized loads are required to saturate the memory bus.

\textbf{Finding RC1: Triton-generated convolution code fails to issue vectorized (128-bit) load instructions,
preventing asynchronous copy (\texttt{cp.async}) from being applied and collapsing the software pipeline.}

% REVISION-NOTE-START id=4134A-W3 severity=Major
\textcolor{red}{\textbf{[REVISION \#4134A-W3, Major]}
Concern (\textit{R1}): RC1 should cite the NCU counter confirming it.
Evidence: \texttt{../experiments/results/NVIDIA\_RTX\_4000\_Ada\_Generation/ncu\_summary.csv} (\texttt{ldg\_per\_request}).
Fix: Add: ``NCU shows \texttt{ldg\_per\_request < 32} bits for Triton conv2d, vs \texttt{$\geq$ 128} bits for Triton matmul --- confirming the absent vectorized load.''}
% REVISION-NOTE-END id=4134A-W3

\phantomsection\label{sec:analysis:autotune}
\paragraph{RC2: Auto-Tuning Search Space Mismatch.}
Triton's auto-tuner sweeps a programmer-specified list of configurations
$\\ \{(\texttt{BLOCK\_M}, \texttt{BLOCK\_N}, \texttt{BLOCK\_K}, \texttt{num\_stages}, \texttt{num\_warps})\}.\\ $ For GEMM, this 5-dimensional space is well-understood:
large square tiles ($128 \times 128$) with 4--8 pipeline stages
are optimal across most A100 problem shapes~\cite{tillet2019triton}.
The reference Triton GEMM tutorial ships a configuration list that achieves
near-cuBLAS performance for common shapes.

This is further evidenced by the large-matmul result in \cref{sec:eval:gemm}:
Triton achieves only 31.7\% of cuBLAS throughput on a $16384^2$ FP16 matrix multiplication
(361.9~ms vs.\ PyTorch's 114.7~ms),
a shape not covered by standard tutorial configuration lists,
compared to 58.2\% on a $128\times2048^2$ batched GEMM
where the auto-tuner has well-populated configurations.
The convolution gap (34.9\% for the $32\times256\times128\times128$, $3\times3$ configuration)
follows the same pattern at an even larger scale.

For the $16384^2$ matmul, the combined A, B, C footprint (${\approx}1.6$~GB) vastly exceeds the 48~MB L2 cache; cuBLAS employs hardware-specific cache-blocking (via \texttt{cuBLASLt}'s plan selector) while Triton's generic \texttt{tl.dot} compilation lacks equivalent heuristics, producing the $3.2\times$ latency penalty observed ($114.7$~ms vs.\ $361.9$~ms).

Convolution adds filter-size degrees of freedom: $1\times1$ convolutions behave like GEMM (large tiles preferred), while $3\times3$ require smaller $K$-dimension tiles to keep shared memory within 48~KB per SM.
Default configuration lists do not cover $5\times5$ and $7\times7$ filters, leaving substantial performance on the table.
TileLang additionally requires \texttt{num\_stages} to be declared ahead of time; for large filters the optimal value depends on register pressure from spatial accumulation---a dependency current scheduling heuristics do not model.

\textbf{Finding RC2: The static tile-configuration search space in Triton's \texttt{@triton.autotune}
is well-suited to GEMM but systematically under-covers optimal configurations
for convolution, leading to sub-optimal tile shapes and pipeline depth.}

% REVISION-NOTE-START id=4134A-W3 severity=Major
\textcolor{red}{\textbf{[REVISION \#4134A-W3+4134A-W7, Major]}
Concern (\textit{R1}): RC2 should distinguish RC2a (configs absent from default grid; e.g.\ matmul L2-swizzle) from RC2b (search budget exhausted before convergence).
Evidence: \texttt{../AKO4ALL/results/optimized/matmul\_triton.py:10-26}, \texttt{../experiments/results/NVIDIA\_RTX\_4000\_Ada\_Generation/autotune\_matmul.csv}.
Fix: Split RC2 into RC2a / RC2b with concrete examples and cite \texttt{autotune\_matmul.csv} for matmul (RC2a) and conv2d (RC2b).}
% REVISION-NOTE-END id=4134A-W3

\phantomsection\label{sec:analysis:regpressure}
\paragraph{RC3: Register Pressure and Occupancy Reduction.}
Each streaming multiprocessor on A100 has 65,536 32-bit registers per block.
A block of 128 threads with 64 registers per thread occupies the full register file,
preventing a second block from being resident and eliminating pipeline interleaving.
For large-filter convolutions, maintaining partial sums for the
$K_H \times K_W$ reduction --- plus input tile and output tile registers ---
pushes per-thread register usage well above 64 registers.
Triton's \texttt{num\_warps} controls thread count per block but does not directly constrain per-thread register usage, leaving register allocation to the compiler.

\textbf{Finding RC3: Convolution kernels with $5 \times 5$ and larger filters
exceed the per-thread register budget, causing spill to local memory
and reducing SM occupancy below the level required for latency hiding.}

% REVISION-NOTE-START id=4134A-W3 severity=Major
\textcolor{red}{\textbf{[REVISION \#4134A-W3+4134A-W13+GAP-2, Major]}
Concern (\textit{R1}): The submitted-PDF RC3 attributes register-spill to conv K$\geq$5, but NCU shows conv \texttt{n\_spills = 0}; spill is TileLang-LayerNorm-specific (254 regs/thread, 51.5\,GB local-load), and L336 cites A100 register file size in an Ada-only study.
Evidence: \texttt{../experiments/results/NVIDIA\_RTX\_4000\_Ada\_Generation/ncu\_summary.csv}.
Fix: Relabel RC3 as ``RC3: TileLang LayerNorm Register Spill (Not a Convolution Problem),'' move conv discussion into RC1/RC2, document layer\_norm spill explicitly (254 regs, 51.5\,GB, 16.5\% occupancy), and change the L336 A100 citation to Ada (sm\_89) or note ``Ada and A100 share the 65{,}536 32-bit regs/SM budget.''}
% REVISION-NOTE-END id=4134A-W3


\phantomsection\label{sec:analysis:algo}
\paragraph{RC4: Absence of Algorithm-Level Diversity.}
cuDNN's runtime selector chooses Winograd for $3 \times 3$ filters
when arithmetic reduction is the bottleneck.
Winograd convolution reduces the multiply-accumulate count by a factor of
approximately $2.25\times$ for $3 \times 3$ (Winograd $F(2,3)$) filters,
at the cost of additional data transformation passes.
Neither Triton nor TileLang expose the Winograd transformation
as a schedulable primitive,
so they cannot exploit this arithmetic reduction.
This contributes to the gap specifically on $3 \times 3$ filters
when the computation is arithmetic-bound.

\textbf{Finding RC4: Both Triton and TileLang implement only the im2col-GEMM lowering for convolution,
providing no access to Winograd or FFT algorithms that cuDNN uses for specific filter sizes.}

% REVISION-NOTE-START id=4134A-W3 severity=Major
\textcolor{red}{\textbf{[REVISION \#4134A-W3+GAP-3, Major]}
Concern (\textit{R1}): RC4 currently presented as a major remaining gap; Winograd isolation shows $\approx$2.28\%.
Evidence: \texttt{../experiments/results/NVIDIA\_RTX\_4000\_Ada\_Generation/winograd\_isolation.csv}.
Fix: Soften the finding to: ``Of the residual gap on 3$\times$3 conv, $\approx$2--3\% is attributable to absent Winograd selection; the remainder reflects cuDNN's general implicit-GEMM tuning.''}
% REVISION-NOTE-END id=4134A-W3

\subsection{Summary of Root Causes}
\label{sec:analysis:summary}

\Cref{tab:rootcauses} maps each root cause to the kernel category it primarily affects,
its estimated contribution to the throughput gap, and the Nsight Compute counter
most diagnostic for its detection.

\begin{table*}[t]
  \centering
  \caption{Root-cause summary. Measured contribution is the latency speedup recovered by the corresponding mitigation (\cref{sec:mitigation}); ``---'' indicates no direct mitigation experiment conducted. RC0a and RC0b are TileLang-specific; RC1--RC4 affect both DSLs.}
  \label{tab:rootcauses}
  \begin{tabular}{p{5cm}p{3.5cm}p{2.2cm}p{4.5cm}}
    \toprule
    Root cause                     & Affected kernels              & Contribution & Diagnostic counter              \\
    \midrule
    RC0a: \texttt{T.serial} instead of \texttt{T.reduce} & All TileLang reductions, norm & $208\times$ (LayerNorm) & Thread sync count, warp stall   \\
    RC0b: Absent vectorized loads / dtype cast            & TileLang norm, conv           & $3.4\times$ (LayerNorm) & \texttt{l1tex\_\_t\_bytes}      \\
    RC1: Strided access, scalar LD  & Conv2d ($K > 1$), Triton      & $2.57\times$ with RC2   & Vectorized-load fraction        \\
    RC2a: Auto-tuning mismatch      & Matmul, Conv2d, Depthwise     & $1.66\times$ (Matmul)   & TFLOPS vs.\ swept configs       \\
    RC2b: L2 cache residency        & Matmul $\geq 16384^2$         & ---                     & L2 hit rate, DRAM throughput    \\
    RC3: Register pressure          & Conv2d ($K \geq 5$)           & ---                     & \texttt{sm\_\_register\_spill} \\
    RC4: No Winograd                & Conv2d $3\times3$             & ---                     & SM utilization delta            \\
    \bottomrule
  \end{tabular}
\end{table*}

% REVISION-NOTE-START id=4134A-W1 severity=Major
\textcolor{red}{\textbf{[REVISION \#4134A-W1+GAP-1, Major]}
Concern (\textit{R1}): The RC0a row's ``Diagnostic counter'' cell reads ``Thread sync count, warp stall,'' implying barrier synchronization is the dominant stall. NCU shows the dominant stall is \texttt{warp\_stall\_long\_scoreboard} (memory latency, $\approx$\,105 cycles), with \texttt{warp\_stall\_barrier} $\approx$ 0.
Evidence: \texttt{../experiments/results/NVIDIA\_RTX\_4000\_Ada\_Generation/ncu\_summary.csv}, \texttt{NCU\_FINDINGS.md}.
Fix: Replace the RC0a ``Diagnostic counter'' cell with \texttt{warp\_stall\_long\_scoreboard} (and note ``barrier $\approx$ 0''). This is the same memory-latency framing applied throughout \S6 above.}
% REVISION-NOTE-END id=4134A-W1

% REVISION-NOTE-START id=4134A-W3 severity=Major
\textcolor{red}{\textbf{[REVISION \#4134A-W3+GAP-2, Major]}
Concern (\textit{R1}): The RC3 row lists ``Affected kernels'' as ``Conv2d ($K \geq 5$)''. NCU shows conv \texttt{n\_spills}\,=\,0; the spill is in TileLang LayerNorm (254 regs/thread, 51.5\,GB local-load, 16.5\% occupancy).
Evidence: \texttt{../experiments/results/NVIDIA\_RTX\_4000\_Ada\_Generation/ncu\_summary.csv}.
Fix: Change the RC3 row's ``Affected kernels'' cell from ``Conv2d ($K \geq 5$)'' to ``TileLang LayerNorm (large)''; add a measured contribution column entry derived from the unspilled-vs-spilled \texttt{sm\_\_register\_spill} ratio for the kernel.}
% REVISION-NOTE-END id=4134A-W3

% REVISION-NOTE-START id=4134A-W3 severity=Major
\textcolor{red}{\textbf{[REVISION \#4134A-W3, Major]}
Concern (\textit{R1}): Counter values in the ``measured contribution'' column should ground every row.
Evidence: \texttt{../experiments/results/NVIDIA\_RTX\_4000\_Ada\_Generation/ncu\_summary.csv}.
Fix: Add a footnote pointing readers to the consolidated \texttt{ncu\_summary.csv} and to per-RC subsections, and update the RC0b row to cite its separate contribution rather than rolling it into RC0.}
% REVISION-NOTE-END id=4134A-W3
